\documentclass[12pt]{article}
\usepackage{bbm}
\usepackage{fancyhdr}
\usepackage{titling}
\usepackage{enumitem}
\usepackage{amsmath}
\usepackage{dsfont}
\usepackage{amsfonts}
\usepackage{amssymb}
\usepackage{amsthm}
\usepackage{geometry}
\usepackage{tikz}
\usepackage{algorithm}
\usepackage{algpseudocode}
\geometry{verbose,tmargin=1in,bmargin=1in,lmargin=1in,rmargin=1in}
\renewcommand{\vec}[1]{\mathbf{#1}}
\newcommand{\dec}{\text{dec}}
\newcommand{\divi}{\mathbin{\textnormal{div}}}
\title{Homework 1}
\author{Satvik Prasad}
\date{}
\newtheorem{lemma}{Lemma}[section]
\fancypagestyle{plain}{%
  \fancyhf{}
  \fancyhead[L]{\theauthor}
  \fancyhead[R]{\thedate}
}
\pagestyle{plain}
\makeatletter
\def\@maketitle{%
  \newpage
  \null
  \vskip 1em%
  \begin{center}%
    \let \footnote \thanks
    {\LARGE \@title \par}%
    \vskip 1em%
  \end{center}%
  \par
\vskip 1em}
\makeatother
\usepackage{cmbright}

\newcommand{\mat}[1]{\mathbf{#1}}

% ---------- Probability theory ----------
% Core operators (optional argument-free; use as \E{X}, \Var{X}, \Cov{X}{Y})
\DeclareMathOperator{\EOp}{\mathbb{E}}
\DeclareMathOperator{\VarOp}{Var}
\DeclareMathOperator{\CovOp}{Cov}
\DeclareMathOperator{\CorrOp}{Corr}

\renewcommand{\P}[1]{\mathbb{P}\left(#1\right)}        % \P{A}
\newcommand{\E}[1]{\EOp\left[#1\right]}                % \E{X}
\newcommand{\Var}[1]{\VarOp\left(#1\right)}            % \Var{X}
\newcommand{\Cov}[2]{\CovOp\left(#1,\,#2\right)}       % \Cov{X}{Y}
\newcommand{\Corr}[2]{\CorrOp\left(#1,\,#2\right)}     % \Corr{X}{Y}
\newcommand{\SD}[1]{\operatorname{SD}\left(#1\right)}  % \SD{X}

% Subscripted versions, e.g. \Esub{\theta}{X}, \Psub{\theta}{A}
\newcommand{\Esub}[2]{\EOp_{#1}\left[#2\right]}
\newcommand{\Psub}[2]{\mathbb{P}_{#1}\left(#2\right)}

% Conditioning, indicators, independence
\newcommand{\given}{\,\middle|\,}                      % \E{X \given Y}
\newcommand{\ind}[1]{\mathbbm{1}\left\{#1\right\}}     % \ind{X > 0}
\newcommand{\indep}{\perp\!\!\!\perp}                  % X \indep Y
\newcommand{\iid}{\overset{\text{iid}}{\sim}}          % X_i \iid \Normal(0,1)

% Convergence
\newcommand{\convp}{\xrightarrow{\;p\;}}               % in probability
\newcommand{\convd}{\xrightarrow{\;d\;}}               % in distribution
\newcommand{\convas}{\xrightarrow{\text{a.s.}}}        % almost surely
\newcommand{\convLp}[1]{\xrightarrow{L^{#1}}}          % in L^p

% Probability spaces
\newcommand{\F}{\mathcal{F}}                           % sigma-algebra
\newcommand{\B}{\mathcal{B}}                           % Borel sigma-algebra
\newcommand{\pspace}{(\Omega, \F, \mathbb{P})}

% Common distributions
\newcommand{\Normal}{\mathcal{N}}
\newcommand{\Bern}{\operatorname{Bernoulli}}
\newcommand{\Bin}{\operatorname{Binomial}}
\newcommand{\Pois}{\operatorname{Poisson}}
\newcommand{\Unif}{\operatorname{Uniform}}
\newcommand{\Exp}{\operatorname{Exp}}
\newcommand{\Geom}{\operatorname{Geometric}}
\newcommand{\Gam}{\operatorname{Gamma}}
\newcommand{\Beta}{\operatorname{Beta}}

% Number sets
\newcommand{\R}{\mathbb{R}}
\newcommand{\N}{\mathbb{N}}
\newcommand{\Z}{\mathbb{Z}}

\begin{document}
\begin{enumerate}
    \item \begin{enumerate}
    \item Assume we have some uniform random variable $U \sim \text{Uniform}(0, 1),$ and some CDF $F$ such that $X = F^{-1}(U)$ is a RV.

    Then, 
    $$
    F_X(k) = \mathbb P(X \leq k) = \mathbb P(F^{-1}(U) \leq k).
    $$
    Since $F^{-1}$ is invertible and monotonically increasing, 
    $$
    \{F^{-1}(U) \leq k\} = \{U \leq F(k)\},
    $$
    meaning we have
    $$
    F_X(k) = F(k),
    $$
    since $F(k) \in [0, 1].$

    \item Assume we have random variable $X$ with CDF $F.$ Then, $Y = F(X)$ has CDF
    $$
    F_y(t) = \mathbb P(F(X) \leq t) = \mathbb P(X \leq F^{-1}(t)) = F(F^{-1}(t)) = t
    $$
    when $0 \leq t \leq 1,$ and $1$ if $t > 1$, and $0$ otherwise.
\end{enumerate}
\item \begin{enumerate}
        \item Assume we have $X \sim \mathcal N(0, 1).$ Let $Y = \sigma X + \mu.$ Then, our CDF satisfies
        $$
        F_Y(x) = \mathbb P(\sigma X+ \mu \leq x) = \mathbb P\left(X \leq \dfrac{x-\mu}{\sigma}\right).
        $$
        Thus, $F_Y(x) = F_X\left(\dfrac{x-\mu}{\sigma}\right),$ and differentiating yields
        $$
        f_Y(x) = \dfrac{1}{\sigma} f_X((x - \mu)/\sigma) = \dfrac{1}{\sigma \sqrt{2\pi}} e^{-\left(\frac{x - \mu}{\sigma}\right)^2},
        $$
        meaning $\sigma X + \mu \sim \mathcal N(\mu, \sigma).$ 

        Similarly, assume we have $Y \sim \mathcal N(\mu, \sigma^2),$ and $X = \dfrac{Y-\mu}{\sigma}.$ Then, the CDF satisfies
        $$
        F_X(x) = \mathbb P\left(\dfrac{Y-\mu}{\sigma} \leq x\right) = \mathbb P(Y \leq \sigma x+ \mu) = F_Y(\sigma x + \mu).
        $$
        Differentiating, we get 
        $$
            f_X(x)= \sigma \dfrac{1}{\sigma \sqrt{2\pi}} e^{-x^2} = \dfrac{1}{\sqrt{2\pi}} e^{-x^2},
        $$
        meaning $X \sim \mathcal N(0, 1)$ as required.

        \item Assume we have $X \sim \text{Exp}(1),$ and $Y = X/\lambda$ for $\lambda > 0.$ Then, we get
        $$
        F_Y(x) = \mathbb P(X \leq \lambda x) = F_X(\lambda x) = 1-e^{-\lambda x},
        $$
        meaning $f_Y(x) = \lambda e^{-\lambda x},$ i.e., $Y \sim \text{Exp}(\lambda).$

        \item Consider $Y = (a-b)U + b.$ Then, $$
        F_Y(x) = F_{U}((b - x) / (a - b)),$$
        for all $a \leq x \leq b,$
        and differentiating yields $$
        f_Y(x)  = \dfrac{1}{a - b}f_U((b - x)/ (a-b)) = \begin{cases}
            \dfrac{1}{a-b} & a \leq x \leq b \\
            0  & \text{otherwise},
        \end{cases}
        $$
        meaning $Y \sim \text{Uniform}(a, b).$
    \end{enumerate}

    \item \begin{enumerate}
        \item We have $$\mathbb P(A \leq x) = \mathbb P\left(X \leq \dfrac{x-5}{2}\right),$$
        meaning
        $$
        F_A(x) = 1-e^{-\lambda(x-5)/2},
        $$
        and thus 
        $$
        f_A(x) = \frac{\lambda}{2} e^{-\lambda(x-5)/2}
        $$
        for $(x - 5)/ 2 \geq 0,$ i.e., $x \geq 5.$ Otherwise, the density is $0.$

        \item We have $$\mathbb P(B \leq x) = \mathbb P(X \leq \ln x)$$ since $e^x$ is monotonically increasing, yielding
        $$
        F_B(x) = 1- e^{-\lambda \ln x} = 1-x^{-\lambda}.
        $$
        Thus, we have density
        $$
        f_B(x) = \lambda x^{-(\lambda + 1)} = \dfrac{\lambda}{x^{\lambda + 1}}
        $$ whenever $\ln x \geq 0,$ or $x \geq 1.$ On $(-\infty, 1],$ the density is $0.$

        \item We have $$\mathbb P(C \leq x) = \mathbb P\left(1 + X \geq \dfrac{1}x\right) = \mathbb P\left(X \geq \dfrac{1}{x} - 1\right),$$ 
        meaning
        \begin{align*}
        F_C(x) &= 1 - F_X(1/x - 1) = 1 - (1 - e^{-\lambda(1/x - 1)}) \\
            &= e^{\lambda -\lambda/x } ,
        \end{align*}
        yielding density
        $$
            f_C(x) = \dfrac{\lambda}{x^2} e^{\lambda(1- 1/x)}
        $$
        when $0 < x \leq 1.$ Outside this domain, density is $0.$

        \item Since $X \geq 0,$ $0 < (1 + X)^{-2} \leq 1.$ We have
        $$
        \mathbb P((1+X)^{-2} \leq x) = \mathbb P\left((1+ X)^2 \geq \dfrac{1}{x}\right).
        $$
        Since $1 + X \geq 0,$ we have
        \begin{align*}
        \mathbb P((1+X)^2 \geq 1/x) &= \mathbb P\left(X \geq \dfrac{1}{\sqrt{x}} - 1\right) = 1 - F_X\left(\dfrac{1}{\sqrt x} - 1\right) \\
        &= e^{-\lambda(1/\sqrt x - 1)}
        \end{align*}
        on $(0, 1].$ Hence, we have density
        $$
            f_D(x) = \dfrac{\lambda}{2x\sqrt{x}} e^{-\lambda(1/\sqrt x -1 )}
        $$
        on $(0, 1]$ and $0$ elsewhere.
    \end{enumerate}

    \item For $X,Y, Z$ independent, we need $f(x, y, z) = f_X(x) f_Y(y) f_Z(z).$ Let us compute marginals. For $X,$
    \begin{align*}
    f_X(x) = \int_{0}^{1} \int_{0}^{1} f(x, y, z) \, dy dz &=  \int_{0}^{1} \int_{0}^{1} 8xyz \, dy \, dz \\
    &= x \int_{0}^{1} 4y \, dz \\
    &= 2x .
    \end{align*}
    By symmetry, $f_Y(y) = 2y, $ and $f_Z(z) = 2z.$ Hence, 
    $$
    f_X(x) f_Y(y) f_Z(z) = 8xyz = f(x, y, z),
    $$
    meaning $X, Y, Z$ are independent.

    \item $X, Y$ are independent, meaning they have joint density $f(x, y ) = f_X(x) f_Y(y).$ We want to find the density of $U = XY,$ so consider the change of variables for $U = XY, V = X$:
    $$
    u(x, y) = xy, \quad v(x, y) =x .
    $$
    Then, we have change of variables map $T : \mathbb R^{2} \to \mathbb R^{2}$ defined by $T(x, y) = (xy, x),$ which is injective when $x \neq  0.$ Hence, we have the inverse
    $T^{-1}: \mathbb R^2 \setminus \{(u, v) : v = 0\} \to \mathbb R^2 \setminus \{(x, y) : x = 0\}$
    $$
        T^{-1}(u, v) = ( v, u / v),
    $$
    with Jacobian
    $$
    J = \left|\dfrac{\partial (x, y)}{\partial (u, v)}\right| = \begin{vmatrix}
    0 & 1 \\
    \dfrac{1}{v} & -\dfrac{u}{v^2}
    \end{vmatrix}.
    $$
    This means we have joint density
    $$
    g(u, v) = \dfrac{f(v, u/v)}{|v|}
    $$
    when $v \neq 0,$ (which has measure $0$). Computing marginal for $U$ yields $$
        f_U(u) = \int_{-\infty}^{\infty} f(x, u/x) \dfrac{1}{|x|} \, dx = \int_{-\infty}^{\infty} f_X(x) f_Y\left(\dfrac{u}{x}\right) \dfrac{1}{|x|} \, dx
    $$
    as required.

    Similarly, for $V = X/Y$ let us construct the change of variables
    $U = X, V = X/ Y,$ with map $T: \mathbb R^2 \setminus \{(x, y) : y  = 0\} \to \mathbb R^{2} \setminus \{(x, y) : y = 0\}$ defined by 
    $$
    T(x, y) = (x, x/y),
    $$
    meaning the inverse satisfies $T^{-1} : D \to D$
    $$
    T^{-1}(u, v) = (u, u/v),
    $$
    where $D = \mathbb R^{2} \setminus \{(x,y) : y \neq 0\}.$ This has Jacobian 
    $$
        J = \begin{vmatrix}
            1 & 0 \\
            \dfrac{1}{v} & -\dfrac{u}{v^2}
        \end{vmatrix}  = -u/v^2.
    $$
    Hence, $U, V$ have joint density
    $$
    g(u, v) =  \dfrac{f(u, u/v) |u|}{v^2},
    $$
    allowing us to compute the marginal 
    as
    $$
    f_V(v) = \int_{-\infty}^{\infty} \dfrac{f_X(x) f_Y(x/v) |x|}{v^2} \, dx.
    $$
    Letting $x = vy,$ we get $\,dx = v \, dy,$
    meaning
    \begin{align*}
    f_V(v) &= \dfrac{v}{|v|}\int_{-\infty}^{\infty} f_X(vy) f_Y(y) |vy| \dfrac{ v }{v^2}\, dy \\
    &=  \int_{-\infty}^{\infty} f_X(vy) f_Y(y) |y| \, dy
    \end{align*}
    as required.

    \item Notice that the event $\{\min\{X, Y\} \leq x\}$ decomposes into 
    $$
        \{X \leq x \cup Y \leq x\}.
    $$
    To make these events disjoint, take
    $$
    \{(X \leq x \cap Y > x) \cup (Y \leq x)\}.
    $$
    Then, by addition under disjoint union, we have
    $$
    \mathbb P(\min\{X, Y\} \leq x) = \mathbb P(X \leq x \cap Y > x) + F_Y(x).
    $$
    Since $(\{X \leq x\}\cap \{Y > x\}) \cup (\{X \leq x\} \cap \{Y \leq x\}) = \{X \leq x\},$ we know
    $$
    \mathbb P(\{X \leq x\} \cap \{Y > x\}) = F_X(x) - \mathbb P(\{X \leq x\} \cap \{Y \leq x\}).
    $$

    Since we have $X, Y$ independent, 
    $$
    \mathbb P(\{X \leq x\} \cap \{Y > x\}) = F_X(x) -F_X(x) F_Y(y),
    $$
    meaning the CDF satisfies
    $$
    \mathbb P(\min\{X, Y\} \leq x) = F_X(x) + F_Y(x)- F_X(x)F_Y(x).
    $$
    Thus, our CDF is $$
    \mathbb P(\min\{X, Y\} \leq x) = 1  - e^{-(\lambda + \mu)x}
    $$
    for $x \geq 0,$ and $0$ otherwise.
    Differentiating yields probability density
    $$
    f(x) = (\lambda + \mu) e^{-(\lambda + \mu)x}
    $$
    for $x \geq 0,$ and $0$ otherwise, meaning $\min\{X, Y\} \sim \text{Exp}(\lambda + \mu)$ as required.
\end{enumerate}
\end{document}